\begin{resumo}

A correção de provas discursivas é uma das tarefas mais onerosas da prática docente, exigindo tempo significativo e estando sujeita à subjetividade do avaliador. Este trabalho apresenta o desenvolvimento e a avaliação de um sistema de correção automatizada de provas discursivas baseado em múltiplos agentes de Inteligência Artificial, com suporte à Recuperação Aumentada por Geração (RAG) e mecanismo de arbitragem por divergência. O sistema utiliza dois agentes corretores independentes que avaliam cada resposta com base em uma rubrica e no material didático da disciplina, acionando condicionalmente um terceiro agente (árbitro) quando a divergência entre as notas ultrapassa um limiar configurável --- desenho inspirado no processo oficial de correção de redações do ENEM. A arquitetura emprega FastAPI para o \textit{backend}, LangGraph para orquestração do \textit{pipeline}, LangChain para integração com modelos de linguagem (Gemini 2.5 Flash) e ChromaDB para armazenamento vetorial. A avaliação foi conduzida por meio de um estudo de caso controlado com três questões discursivas de Algoritmos e Estrutura de Dados, quatro níveis de qualidade de resposta (excelente, adequada, fraca e fora do tema) e duas condições experimentais (com e sem RAG). Os resultados indicam que o sistema completou o fluxo \textit{end-to-end} sem falhas (100\% de completude estrutural), apresentou alta estabilidade nas notas (variação máxima de 0,21 ponto em três repetições) e demonstrou que o RAG contribui positivamente para a avaliação, especialmente em respostas de qualidade intermediária. O protótipo preserva a autoridade docente, disponibilizando os resultados para revisão, ajuste e aprovação antes da consolidação final.

\textbf{Palavras-chave:} Correção Automatizada. Inteligência Artificial. Modelos de Linguagem (LLM). Recuperação Aumentada por Geração (RAG). Sistemas Multiagentes. Avaliação Educacional.

\end{resumo}
